\documentclass{article}
\usepackage[utf8]{inputenc}
\usepackage[spanish, es-noshorthands]{babel}
\usepackage{tikz}
\usetikzlibrary{arrows.meta, positioning, calc}
\usetikzlibrary{babel} 

\usepackage[margin=0.5cm]{geometry}

\begin{document}
\pagestyle{empty} 
\centering

\begin{tikzpicture}[
    >=Stealth,
    font=\sffamily\small,
    % Estilo para las cajas de componentes
    box/.style={
        draw=black, fill=gray!4, rounded corners=6pt, line width=0.6pt,
        align=center, inner sep=8pt, text width=4cm, minimum height=1cm
    },
    arrow/.style={->, line width=0.7pt, draw=black}
]

% --- CAPA 1: Orígenes ---
\node[box] (app) at (-3, 6) {Application Services};
\node[box] (kafka) at (3, 6) {Kafka Components};

% --- CAPA 2: Intermedios (Alineados perfectamente debajo) ---
\node[box] (verif) at (-3, 4) {Verification\\Framework};
\node[box] (deploy) at (3, 4) {Deployment\\Infrastructure};

% --- CAPA 3 y 4: Destinos Centrales ---
\node[box, text width=6cm] (diagnostic) at (0, 1.5) {Central Diagnostic Collection};
\node[box, text width=7cm] (users) at (0, 0) {Operators / Developers / Administrators};

% --- CONEXIONES VERTICALES DIRECTAS ---
\draw[arrow] (app.south) -- (verif.north);
\draw[arrow] (kafka.south) -- (deploy.north);

% --- EL ACOPLADOR EN T SIMÉTRICO ---
% Calculamos el punto de unión central a una altura intermedia (Y = 2.8)
\coordinate (union) at (0, 2.8);

% Forzamos a que las líneas salgan del SUR de las cajas, bajen rectas y giren al centro
\draw[line width=0.7pt, draw=black] (verif.south) |- (union);
\draw[line width=0.7pt, draw=black] (deploy.south) |- (union);

% Del punto de unión central sale la flecha vertical hacia el bloque de diagnóstico
\draw[arrow] (union) -- (diagnostic.north);

% Conexión final inferior
\draw[arrow] (diagnostic.south) -- (users.north);

\end{tikzpicture}

\end{document}
